\documentclass[11pt,a4paper]{article}
\usepackage[T1]{fontenc}
\usepackage[utf8]{inputenc}
\usepackage[french]{babel}
\frenchsetup{SmallCapsFigTabCaptions=false}
\usepackage{lmodern,microtype}
\usepackage[left=22mm,right=22mm,top=23mm,bottom=24mm]{geometry}
\usepackage{amsmath}
\usepackage{graphicx,booktabs,array,tabularx,float}
\usepackage{xcolor,tikz}
\usepackage{enumitem,titlesec}
\usepackage{hyperref}
\hypersetup{colorlinks=true,linkcolor=black,citecolor=black,urlcolor=sfigBlue,pdfauthor={Projet LANCE},pdftitle={Structurer le corpus de scenarios LANCE avant le gel},pdflang={fr-FR}}
\urlstyle{same}
\setlength{\parindent}{1.2em}
\setlength{\parskip}{3pt}
\setlength{\emergencystretch}{2em}
\clubpenalty=10000
\widowpenalty=10000
\displaywidowpenalty=10000
\setlist{nosep,leftmargin=5mm,topsep=3pt}
\def\UrlBreaks{\do\/\do\-\do\_\do\.\do\?\do\&\do\=\do\:\do\,}
\newcolumntype{Y}{>{\raggedright\arraybackslash}X}
\newcommand{\code}[1]{{\small\ttfamily\path{#1}}}
\titleformat{\section}{\Large\bfseries}{\thesection}{0.7em}{}
\titleformat{\subsection}{\large\bfseries}{\thesubsection}{0.7em}{}
\titlespacing*{\section}{0pt}{16pt}{7pt}
\titlespacing*{\subsection}{0pt}{11pt}{5pt}
\pagestyle{plain}
\setcounter{tocdepth}{2}
\usetikzlibrary{arrows.meta,positioning,fit,calc}
\definecolor{sfigInk}{HTML}{1F242C}
\definecolor{sfigMuted}{HTML}{5C626C}
\definecolor{sfigLine}{HTML}{D9DEE5}
\definecolor{sfigBlue}{HTML}{1D5C96}
\definecolor{sfigBlueFill}{HTML}{EFF6FC}
\definecolor{sfigRoseFill}{HTML}{FDF0F4}
\tikzset{
  sfig diagram/.style={>=Latex, node distance=8mm and 9mm, every node/.style={font=\sffamily\footnotesize,text=sfigInk}},
  gbar/.style={draw=sfigLine, fill=sfigBlueFill, rounded corners=2pt, line width=0.55pt, align=center, inner sep=4pt},
  gmil/.style={draw=sfigBlue, fill=white, rounded corners=2pt, line width=0.7pt, align=center, inner sep=4pt}
}

\title{Structurer le corpus de sc\'enarios LANCE\\avant le gel\\[4pt]
\large Grille d\'eriv\'ee, edge cases d\'eployables, sets et s\'equen\c c{}age}
\author{Projet LANCE}
\date{Derni\`ere r\'evision : 2026-10-12 23:05 (Europe/Brussels, UTC+02:00)\\
Statut : propositions relues --- aucun r\'esultat exp\'erimental nouveau}

\begin{document}
\maketitle

\begin{abstract}
Le corpus actuel compte 29 sc\'enarios et 288 vuln\'erabilit\'es attendues, dont 75~\% \`a signature : il est adapt\'e au d\'eveloppement, pas \`a l'\'evaluation comparative vis\'ee \`a AsiaCCS 2027. Ce rapport d\'ecide comment le structurer avant tout gel. Il retient une grille d'\'etiquetage \emph{d\'eriv\'ee} des v\'erit\'es terrain existantes, stock\'ee dans un index annexe r\'eg\'en\'erable sans modifier le catalogue ni les v\'erit\'es terrain ; quatre edge cases retenus --- EC1 et EC2 d\'eployables, EC3 et EC4 \`a valider en ex\'ecution r\'eelle (variation de surface, bruit et leurres, cha\^inage cassable, autorisation presque valide) ---, deux autres explicitement diff\'er\'es apr\`es la soumission ; un set pilote de calibration (S1, S14, S15, S1h) et un set campagne \`a split par famille dimensionn\'e par la variance du pilote ; enfin un s\'equen\c c{}age en trois horizons avec crit\`eres d'arr\^et \`a la d\'ecision du 2026-11-16. Les propositions ont \'et\'e soumises \`a deux contre-revues (m\'ethodologie, faisabilit\'e) dont les corrections sont int\'egr\'ees.
\end{abstract}

\tableofcontents

\section{Objectif, p\'erim\`etre et m\'ethode}

Ce document r\'epond \`a la question : que faut-il structurer, compl\'eter ou geler dans le corpus de sc\'enarios --- et dans quel ordre --- pour que la comparaison entre la politique \`a r\`egles D1 (\code{rules}) et la politique \`a mod\`ele A1 (\code{llm}) soit recevable \`a AsiaCCS 2027 (soumission le 2026-12-11, d\'ecision interm\'ediaire le 2026-11-16) ? Il prolonge l'audit de couverture du dossier d'objectifs (\code{objectifs/asiaccs-2027/audit-couverture.md} : 215 failles sur 288 \`a signature, strate logique r\'eduite au seul S15, axes variation/bruit/protocole non encod\'es, cha\^inage confondu avec la taille des sc\'enarios) et la proposition de recentrage de ce dossier (\code{proposition-conditions.tex}).

M\'ethode : trois propositions parall\`eles (structuration, edge cases, s\'equen\c c{}age), deux contre-revues ind\'ependantes (m\'ethodologie : contamination et \'etiquetage arrangeant ; faisabilit\'e : effort r\'eel sur l'infrastructure Ansible/Proxmox et d\'elai), puis synth\`ese. Les corrections des relecteurs sont int\'egr\'ees au corps du texte. Statut : propositions --- aucun run, aucun d\'eploiement et aucune g\'en\'eration n'ont \'et\'e ex\'ecut\'es pour ce rapport ; les efforts (en jours-personne, j-h) sont des estimations \`a confirmer. D\'efinitions minimales : une \emph{famille} est un m\'ecanisme de faille (pack) ; une \emph{instance} est une variante d\'eploy\'ee d'une famille ; un \emph{edge case} est une situation con\c c{}ue pour s\'eparer les comportements (r\'eviser, s'abstenir, distinguer), pas seulement les scores.

\section{Grille d'\'etiquetage retenue : d\'eriver, ne pas r\'e\'ecrire}
\label{sec:grille}

Premier constat v\'erifi\'e : le catalogue (\code{benchmarks/catalog.yaml}, charg\'e par \code{src/benchmark/catalog.py}) impose un sch\'ema strict de cl\'es (\code{id}, \code{label}, \code{split}, \code{profile}) et rejette toute cl\'e inconnue. Ajouter des champs d'\'etiquetage au catalogue casserait donc son chargement. D\'ecision : \textbf{catalogue et v\'erit\'es terrain inchang\'es} ; la grille est une vue \emph{calcul\'ee}.

Les champs existent d\'ej\`a dans les v\'erit\'es terrain : par faille (\code{category}, \code{hop_depth}, \code{severity}, \code{template_key}, \code{security_profile}), par sc\'enario (\code{difficulty}, \code{controls}, \code{attack_paths}, \code{scoring}). Le recalcul sur les 31 fichiers (29 sc\'enarios, 2 durcis) confirme l'audit : 70 misconfigurations, 54 expositions de donn\'ees, 50 absences d'authentification, 41 identifiants par d\'efaut ; profondeur de saut 0 : 166 failles, 1 : 67, 2 : 49, 3 : 6 --- le cha\^inage est port\'e par les gros sc\'enarios de d\'eveloppement (S11--S13), d'o\`u la confusion taille/profondeur \`a d\'ecoupler.

La grille d\'eriv\'ee comporte cinq axes, jamais saisis \`a la main :

\begin{itemize}
\item \textbf{Type} : \code{signature} ou \code{logique}, d\'eduit de la cat\'egorie (logique = autorisation, tenant, JWT, IDOR ; aujourd'hui S15 seul).
\item \textbf{Cha\^inage} : profondeur de saut (\code{hop_depth}), utilis\'ee comme axe propre, d\'ecoupl\'ee de la taille du sc\'enario dans l'analyse.
\item \textbf{Surface} : \code{base}, \code{rotate_ips}, \code{rename_hosts} --- lu depuis le manifeste d'alt\'eration appliqu\'e \`a l'instance, pas depuis la v\'erit\'e terrain.
\item \textbf{Bruit} : \code{none}, \code{decoys}, \code{near_miss}, \code{noise}, \code{evidence_degradee} --- m\^eme source (manifeste d'alt\'eration).
\item \textbf{Protocole} : \code{connu} ou \code{inconnu}, explicite ; aujourd'hui non encod\'e, valeur par d\'efaut \code{connu} (limite document\'ee, section~\ref{sec:limites}).
\end{itemize}

Stockage : un index annexe g\'en\'er\'e et versionn\'e (par exemple \code{benchmarks/labels/derived_index.json}), cl\'e (sc\'enario, faille), r\'eg\'en\'erable depuis les v\'erit\'es terrain et les manifestes d'alt\'eration, avec empreintes des sources. L'unit\'e d'analyse r\'eutilise celle, d\'ej\`a op\'erationnelle, de la campagne (\code{family_id} / \code{instance_id} / \code{trial_id}, agr\'egation par famille dans \code{src/benchmark/policy_campaign.py}). Effort estim\'e : script de d\'erivation et tests de non-r\'egression du chargement du catalogue, 2 \`a 3 j-h ; validation crois\'ee contre l'audit de couverture, 0,5 j-h.

\section{Edge cases retenus, par ordre de d\'eployabilit\'e}
\label{sec:edge}

Six edge cases \'etaient candidats (intermittence, bannière mensong\`ere, identifiants presque valides, faille non confirmable, h\^ote incoh\'erent, leurre plausible). Les contre-revues ont impos\'e un filtre d\'ecisif : une alt\'eration seulement pr\'evue (\code{PREVIEW_ONLY}) n'est pas d\'eployable --- le garde-fou d'ex\'ecution r\'eelle (\code{manual_execution.py}) bloque ces alt\'erations, et un r\'esultat en aper\c c{}u ne vaut pas preuve. Le tableau~\ref{tab:edge} retient quatre cas, dans l'ordre.

\begin{table}[H]
\centering\small
\begin{tabularx}{\textwidth}{@{}p{26mm}p{50mm}Y@{}}
\toprule
\textbf{Cas} & \textbf{G\'en\'eration} & \textbf{V\'erification exig\'ee} \\
\midrule
EC1 Variation de surface & Alt\'erations d\'eployables \code{rotate_ips}, \code{rename_hosts} sur une instance d'une famille & V\'erit\'e terrain inchang\'ee hors adresse/nom ; chargement et d\'eploiement r\'eels OK ; m\^eme scoring \\
EC2 Bruit et leurres & \code{add_decoys}, \code{near_miss_controls}, \code{add_noise}, \code{degrade_evidence} ; instance s\oe ur de EC1 & Failles constantes en nombre et cat\'egorie ; aucun leurre compt\'e comme faille (contr\^ole n\'egatif) \\
EC3 Cha\^inage cassable & Pack \code{f20} (pivots SSH/HTTP, profondeur 0--3) : rompre un seul maillon en gardant la profondeur cible & Chemins d'attaque de la v\'erit\'e terrain coh\'erents, profondeur recalcul\'ee, aucun chemin alternatif r\'esiduel \\
EC4 Autorisation presque valide & Pack \code{f11} et S15 (simulateur d\'eterministe porteur/IDOR/JWT) : variantes \`a port\'ee presque suffisante, tenant voisin & R\'eponse d\'eterministe du simulateur ; axe type = logique ; la classe logique d\'epasse enfin $n=1$ \\
\bottomrule
\end{tabularx}
\caption{Edge cases retenus avant la soumission, par d\'eployabilit\'e d\'ecroissante. Efforts estim\'es : EC1 et EC2, 1 \`a 2 j-h chacun (existant, surtout tests) ; EC3, 2 \`a 3 j-h ; EC4, 2 \`a 4 j-h. V\'erification transverse : chaque cas doit passer le garde-fou d'ex\'ecution r\'eelle, pas seulement l'aper\c c{}u.}
\label{tab:edge}
\end{table}

Diff\'er\'es apr\`es la soumission : \textbf{EC5} (d\'efaillance temporelle : \code{rate_limit}, \code{startup_delay}, \code{service_availability} sont \'ecrits en topologie sans \^etre d\'eploy\'es) et \textbf{EC6} (protocole inconnu : exige un nouveau mode du simulateur HTTP unique, d\'eploy\'e par Ansible sur conteneurs LXC --- le chantier le plus lourd, pour l'axe le moins couvert). Ils sont document\'es comme limites, non combl\'es ; seul un essai born\'e de EC5 est admissible si le pilote finit t\^ot.

\section{Sets tir\'es de la grille}
\label{sec:sets}

\textbf{Pilote (calibration, aucune puissance revendiqu\'ee).} S1 (routine), S14 (faible pr\'evalence), S15 (logique, unique), S1h (durci existant, z\'ero faille) qui \emph{remplace} le contr\^ole sain fictif encore pr\'esent dans le manifeste exemple (\code{ground_truth: null}) : un contr\^ole sans v\'erit\'e terrain ne doit jamais \^etre agr\'eg\'e. Format h\'erit\'e : 12 paires, 24 runs, ordre \`a randomiser avant gel (le manifeste actuel est non gel\'e, configuration vide, contr\^oles en attente, ordre altern\'e non randomis\'e). Effort : gel du manifeste et randomisation, 1 j-h ; ex\'ecution et d\'epouillement, 2 \`a 3 j-h.

\textbf{Campagne (apr\`es la d\'ecision du 2026-11-16 seulement).} Split \textbf{par famille}, jamais par instance, pour \'eviter la fuite entre variantes d'une m\^eme famille ; instances proc\'edurales g\'en\'er\'ees par famille via EC1 et EC2 (plus EC3/EC4 sur les familles concern\'ees), car 29 sc\'enarios seuls ne donnent pas la puissance sur les hypoth\`eses H0--H3. Rappel des hypoth\`eses test\'ees (proposition compagne) : H0 aucun \'ecart au-del\`a de $\delta$ ; H1 \'ecart conditionnel (nouveaut\'e, logique, protocole inconnu) ; H2 r\'esilience diff\'erentielle sous bruit ; H3 co\^ut de l'avantage. Taille fix\'ee par calcul de puissance sur la variance du pilote, jamais \`a l'avance. Crit\`ere pr\'e-enregistr\'e : borne inf\'erieure de l'intervalle de confiance appari\'e \`a 95~\% sup\'erieure \`a $\delta$ (valeur de travail 5 points) apr\`es correction de Holm ; quatre contr\^oles n\'egatifs invalidant le lot entier si l'un \'echoue --- dont les durcis S1h et S4h (second sc\'enario durci, ICS/SCADA, z\'ero faille) : toute \og faille \fg{} rapport\'ee sur un durci invalide le lot. Effort : g\'en\'eration et contr\^oles, 5 \`a 8 j-h, ex\'ecution selon la taille calcul\'ee.

\section{S\'equen\c c{}age en trois horizons}
\label{sec:seq}

\begin{itemize}
\item \textbf{Horizon 1 --- jusqu'au 2026-11-16 (d\'ecision).} Grille d\'eriv\'ee et index annexe ; EC1 et EC2 activ\'es sur le pilote, EC3 et EC4 prototyp\'es ; manifeste pilote gel\'e et 12 paires ex\'ecut\'ees. Effort 8 \`a 12 j-h. D\'ecision : \emph{GO} si la borne pilote est compatible avec $\delta$ et les quatre contr\^oles tiennent ; \emph{REPOSITIONN\'E} si l'effet est r\'eel mais sous $\delta$, ou port\'e par le seul axe signature (contribution recentr\'ee sur le benchmark) ; \emph{NO-GO} si un contr\^ole tombe ou si l'effet dispara\^it apr\`es d\'econfusion taille/profondeur.
\item \textbf{Horizon 2 --- 2026-11-17 au 2026-12-11 (soumission).} Uniquement si GO ou REPOSITIONN\'E : campagne \`a split par famille, r\'edaction des 12 pages ACM, paquet Open Science, anonymisation double aveugle, ORCID. Gel des r\'esultats au plus tard 7 \`a 10 jours avant le d\'ep\^ot ; aucun nouvel edge case apr\`es ce gel. Effort 10 \`a 15 j-h hors temps machine.
\item \textbf{Horizon 3 --- apr\`es le 2026-12-11.} EC5 et EC6, poursuite conditionnelle des phases 2 et 5. Crit\`ere d'entr\'ee : soumission d\'epos\'ee ; crit\`ere d'arr\^et propre : fournisseur r\'eel test\'e en ex\'ecution, ou abandon document\'e. Non budg\'et\'e avant le d\'ep\^ot.
\end{itemize}

La figure~\ref{fig:horizons} r\'esume ce s\'equen\c c{}age.

\begin{figure}[H]
\centering
\begin{tikzpicture}[sfig diagram]
\node[gbar, text width=36mm] (h1) {Horizon 1\\Grille + EC1/EC2\\Pilote 12 paires};
\node[gbar, text width=36mm, right=of h1] (h2) {Horizon 2\\Campagne par famille\\12 pages + artefact};
\node[gbar, text width=36mm, right=of h2] (h3) {Horizon 3\\EC5 / EC6\\Phases 2 et 5};
\node[gmil, text width=36mm, below=of h2] (dec) {2026-11-16\\GO / REPOSITIONN\'E\\NO-GO};
\draw[->, draw=sfigMuted] (h2) -- (h3);
\draw[->, draw=sfigMuted] (h1) |- (dec);
\draw[->, draw=sfigMuted] (dec) -| (h2);
\end{tikzpicture}
\caption{S\'equen\c c{}age propos\'e en trois horizons, articul\'e sur la d\'ecision du 2026-11-16. Calendrier de travail, pas un r\'esultat.}
\label{fig:horizons}
\end{figure}

\section{Ce qu'on ne fait pas avant le 11 d\'ecembre}

Six exclusions, chacune motiv\'ee par la validit\'e ou le calendrier : (i) aucun fournisseur (\emph{provider}) pour les alt\'erations \code{PREVIEW_ONLY} (\code{rotate_ports}, segmentation, disponibilit\'e, d\'elais, limitation, impairment) --- effort incompatible avec cinq semaines avant d\'ecision, preuve en aper\c c{}u irrecevable ; (ii) aucun nouveau mode \og protocole inconnu \fg{} du simulateur ; (iii) aucune modification du sch\'ema du catalogue ni des v\'erit\'es terrain pour y encoder surface, bruit ou protocole --- ces axes vivent dans le manifeste d'alt\'eration et l'index d\'eriv\'e ; (iv) aucune extension du pilote au-del\`a de 12 paires ni campagne anticip\'ee --- sans variance pilote, le calcul de puissance est fictif ; (v) aucun bonus hors chemin critique (phases 2/5, durcissements nouveaux au-del\`a de S1h/S4h, refonte de l'agr\'egation qui fonctionne) ; (vi) aucun usage du contr\^ole sain fictif comme r\'esultat.

\section{Limites et points non r\'esolus}
\label{sec:limites}

Les alt\'erations \code{PREVIEW_ONLY} restent bloqu\'ees en ex\'ecution r\'eelle, sans fournisseur livr\'e avant le 2026-12-11. Protocole inconnu, variation de surface et bruit ne sont pas encod\'es dans les v\'erit\'es terrain ; seuls surface et bruit sont d\'erivables du manifeste d'alt\'eration, l'axe protocole reste une limite document\'ee. La taille de campagne n'est pas fixable avant la variance du pilote. Le manifeste pilote reste \`a geler (configuration vide, contr\^oles en attente, ordre \`a randomiser). EC4 doit d\'emontrer en ex\'ecution r\'eelle que la classe logique d\'epasse $n=1$ ; sinon la condition logique reste \`a un seul repr\'esentant. Enfin, la d\'econfusion taille/profondeur par l'index d\'eriv\'e n'est pas encore valid\'ee sur le pilote.

\section{Conclusion}

Le corpus n'a pas besoin d'\^etre reconstruit, mais lu \`a travers une grille : \'etiqueter en d\'erivant, compl\'eter par quatre edge cases d\'eployables dans l'ordre, tirer les sets de la grille, et geler seulement \`a la veille de mesurer. La prochaine action est le script d'index d\'eriv\'e et le gel du manifeste pilote sur S1, S14, S15, S1h --- pas un sc\'enario de plus.

\end{document}
